\chapter{تبدیل متن به گفتار: مسئله و معماری خط‌لولهٔ \eng{TTS}}

\section{از تحلیل گفتار تا تولید گفتار}

تا این جلسه محور درس \textbf{تحلیل سیگنال گفتار} بود: موج صوتی از پیش در اختیار ما بود و از آن طیف‌نگار، \eng{MFCC}، پروزودی، فرمنت‌ها و ... استخراج می‌کردیم. از این جلسه وارد سامانه‌های اصلی می‌شویم: تبدیل متن به گفتار (\eng{TTS})، سامانه‌های گفت‌وگو و تبدیل گفتار به متن (\eng{STT}). موضوع این جلسه \term{تبدیل متن به گفتار}{Speech Synthesis} است. در \eng{TTS} جهت مسئله معکوس می‌شود:
\[
\text{\lr{Text}} \;\longrightarrow\; \text{\lr{Speech Waveform}}
\]
دو مشکل بنیادی وجود دارد:
\begin{itemize}
  \item \textbf{شکاف حجم اطلاعات:} متن بسیار کم‌حجم است (چند ده نویسه یا واج)، اما با نرخ نمونه‌برداری $f_s = 24\,000$ هرتز هر ثانیه گفتار ۲۴ هزار نمونه دارد.
  \item \textbf{کم‌گویی متن:} متن لحن، تأکید، سرعت، مکث و ویژگی‌های گوینده را مشخص نمی‌کند؛ پس مدل باید این‌ها را از داده و بافت \textbf{استنتاج یا تولید} کند. نگاشت متن به گفتار یک‌به‌چند است.
\end{itemize}


\section{دو رویکرد کلی: خط‌لوله و سرتاسری}

\begin{itemize}
  \item \textbf{خط‌لوله‌ای (\eng{Pipeline-Based}):} مسئله به چند ماژول با وظیفهٔ مشخص تقسیم می‌شود و خروجی هر ماژول ورودی بعدی است. هر بخش را می‌توان جداگانه بررسی، آموزش، اصلاح یا جایگزین کرد.
  \item \textbf{سرتاسری (\eng{End-to-End}):} یک مدل (یا مجموعهٔ یکپارچه‌ای از لایه‌ها) متن را مستقیماً به گفتار می‌رساند و لزوماً ماژول‌های میانی مستقل و قابل مشاهده ندارد.
\end{itemize}

\begin{exambox}{نکتهٔ تأکیدی استاد}
رویکرد \eng{End-to-End} \textbf{لزوماً بهترین گزینه برای همهٔ کاربردها نیست}. گاهی به‌دلیل تأخیر پردازش (\eng{Latency}) یا ویژگی‌های دامنهٔ کاربرد (مثلاً حساسیت مرحلهٔ نرمال‌سازی متن)، ماژول‌های کلاسیک یا خط‌لوله مناسب‌ترند. در این جلسه تمرکز بر خط‌لوله است و رویکردهای سرتاسری در جلسهٔ جدا می‌آیند.
\end{exambox}

\begin{correctionbox}{«سرتاسری» یک طیف است}
برخی مدل‌ها (مثل \eng{Tacotron 2}) فقط تا طیف‌نگار سرتاسری‌اند و وکودر جدا دارند؛ مدل‌هایی مانند \eng{VITS} واج را در یک مدل به موج می‌رسانند، و حتی آن‌ها نرمال‌سازی و \eng{G2P} را بیرون از مدل می‌گذارند. پس مسئله بیشتر این است که مرز ماژول‌ها کجا باشد.
\end{correctionbox}

\section{معماری کلی یک سامانهٔ \eng{TTS}}

سامانهٔ \eng{TTS} دو بخش اصلی دارد:
\begin{description}
  \item[\eng{Front-End} (تحلیل متن، \eng{Text Analysis}).] ورودی هنوز متن است؛ «چه باید گفته شود». خروجی: توالی واج‌ها.
  \item[\eng{Back-End} (تولید گفتار، \eng{Speech Synthesis}).] واج‌ها به نمایش آکوستیکی و سپس به موج تبدیل می‌شوند؛ «چگونه گفته شود».
\end{description}

\begin{figure}[!htbp]
\centering
\resizebox{\textwidth}{!}{%
\begin{tikzpicture}[rtl,
  box/.style={draw, rounded corners=3pt, minimum width=2.55cm, minimum height=1.0cm, align=center, fill=softblue, font=\small},
  data/.style={draw, rounded corners=1pt, minimum width=2.2cm, minimum height=0.8cm, align=center, fill=softyellow, font=\small},
  arr/.style={-{Stealth[length=2.5mm]}, thick},
  small/.style={font=\scriptsize, text=black!65, align=center}
]
% Front-end (right to left)
\node[data] (txt) at (12.2,0) {\rl{متن خام}\\\scriptsize\lr{Text}};
\node[box] (norm) at (9.1,0) {\rl{نرمال‌سازی متن}\\\scriptsize\lr{Text Normalization}};
\node[box] (homo) at (6.0,0) {\rl{رفع ابهام هم‌نویسه}\\\scriptsize\lr{Homograph Disamb.}};
\node[box] (g2p) at (2.9,0) {\rl{تبدیل نویسه به واج}\\\scriptsize\lr{G2P}};
\node[data] (ph) at (0,0) {\rl{توالی واج‌ها}\\\scriptsize\lr{Phonemes}};
\draw[arr] (txt) -- (norm);
\draw[arr] (norm) -- (homo);
\draw[arr] (homo) -- (g2p);
\draw[arr] (g2p) -- (ph);
% Back-end (left to right)
\node[box, fill=softgreen] (am) at (3.1,-2.7) {\rl{مدل آکوستیکی}\\\scriptsize\lr{Acoustic Model}};
\node[data] (mel) at (6.1,-2.7) {\rl{طیف‌نگار مل}\\\scriptsize\lr{Mel-Spectrogram}};
\node[box, fill=softgreen] (voc) at (9.1,-2.7) {\rl{وکودر}\\\scriptsize\lr{Vocoder}};
\node[data] (wav) at (12.2,-2.7) {\rl{موج گفتار}\\\scriptsize\lr{Waveform}};
\draw[arr] (ph.south) |- (am.west);
\draw[arr] (am) -- (mel);
\draw[arr] (mel) -- (voc);
\draw[arr] (voc) -- (wav);
% region labels
\draw[dashed, gray] (-1.4,0.95) rectangle (13.5,-0.95);
\node[anchor=south west, font=\small\bfseries, text=mainblue] at (-1.3,0.95) {\lr{Front-End: Text Analysis}};
\draw[dashed, gray] (1.7,-1.75) rectangle (13.5,-4.15);
\node[anchor=north west, font=\small\bfseries, text=mainblue] at (1.8,-4.2) {\lr{Back-End: Speech Generation}};
% size notes
\node[small] at (-0.9,-1.3) {\rl{$30$--$50$ واج}};
\node[small, anchor=north] at (6.1,-3.3) {\rl{$\approx 80$ قاب بر ثانیه}};
\node[small, anchor=north] at (12.2,-3.3) {\rl{$\approx 24000$ نمونه بر ثانیه}};
\end{tikzpicture}}
\caption{نمودار بلوکی خط‌لولهٔ \eng{TTS}: تحلیل متن (بالا، از راست به چپ) و تولید گفتار (پایین، از چپ به راست). مستطیل‌های زرد نمایش‌های میانی و رنگی‌ها ماژول‌های پردازشی‌اند. (شکل افزوده)}
\end{figure}

نقش ماژول‌ها (مراحل ۱ تا ۳ \eng{Front-End} و ۴ و ۵ \eng{Back-End}):
\begin{enumerate}
  \item \term{نرمال‌سازی متن}{Text Normalization}: تبدیل متن خام به متنی که درست خوانده شود (عدد، تاریخ، ساعت، پول، اختصار).
  \item \term{رفع ابهام هم‌نویسه}{Homograph Disambiguation}: تعیین تلفظ درست کلمات هم‌نویسه از روی بافت.
  \item \term{تبدیل نویسه به واج}{Grapheme-to-Phoneme, G2P}: نمایش نوشتاری $\leftarrow$ نمایش واجی.
  \item \term{مدل آکوستیکی}{Acoustic Model}: واج‌ها $\leftarrow$ نمایش آکوستیکی (مثل طیف‌نگار مل).
  \item \term{وکودر}{Vocoder}: نمایش آکوستیکی $\leftarrow$ موج.
\end{enumerate}

\section{مرحلهٔ اول: نرمال‌سازی متن}

متن واقعی پر از عدد، تاریخ، ساعت، واحد پول و اختصار است که شکل نوشتاری‌شان با شکل گفتاری‌شان فرق دارد. در کلاس جملهٔ نمونه‌ای روی تخته نوشته شد که تقریباً همهٔ این پدیده‌ها را داشت.

\begin{figure}[!htbp]
\centering
\newcommand{\chmk}[2]{\tikz[baseline=(n.base)]{\node[draw=red!70!black,ellipse,inner xsep=1.5pt,inner ysep=0pt,line width=0.6pt](n){#2};\node[font=\tiny\bfseries,text=red!70!black,anchor=south west,inner sep=0.5pt] at (n.north east){\lr{#1}};}}
\fbox{\begin{minipage}{0.9\textwidth}
\medskip
\large\linespread{1.5}\selectfont
در تاریخ \chmk{1}{$1402/10/15$} ساعت \chmk{2}{$16{:}30$}، آن \chmk{3}{مرد} پس از مدت‌ها بیماری در بیمارستان \chmk{4}{مرد} و \chmk{5}{سرم} را از دستش کشیدند. \chmk{6}{خواهرم} گفت «\chmk{7}{سرم} به شدت درد می‌کند». حدود \chmk{8}{$2/5$} \chmk{9}{تن} گندم به ارزش \chmk{10}{۱۲۵} میلیون و \chmk{11}{۵۰۰} هزار \chmk{12}{تومان} با \chmk{13}{۳۰٪} تخفیف فروخته شد.
\medskip
\end{minipage}}

\medskip
\small
\begin{tabular}{@{}clll@{}}
\toprule
\textbf{شمارهٔ نشانه} & \textbf{نوع} & \textbf{شکل نوشتاری} & \textbf{شکل گفتاری (یادداشت استاد)}\\
\midrule
۱ & تاریخ & $1402/10/15$ & پانزدهم دی‌ماه هزار و چهارصد و دو\\
۲ & ساعت & $16{:}30$ & چهار و نیم (بعدازظهر) یا شانزده و سی دقیقه\\
۳--۴ & هم‌نویسه & مرد & \lr{mard} (انسان مذکر) / \lr{mord} (فوت کرد)\\
۵، ۷ & هم‌نویسه & سرم & \lr{sorom} (محلول تزریقی) / \lr{saram} (سرِ من)\\
۸ & عدد اعشاری & $2/5$ & دو و نیم\\
۱۰--۱۱ & عدد & ۱۲۵ و ۵۰۰ & صد و بیست و پنج میلیون و پانصد هزار\\
۱۲--۱۳ & پول، درصد & تومان، ۳۰٪ & تومان (محاوره‌ای: تومن)، سی درصد\\
\bottomrule
\end{tabular}
\caption{جملهٔ نمونهٔ روی تخته و نشانه‌گذاری‌های استاد: تاریخ، ساعت، عددها، واحد پول، درصد، و دو هم‌نویسهٔ «مرد» و «سرم». (بازسازی‌شده از تختهٔ کلاس)}
\end{figure}

\begin{definitionbox}{هدف نرمال‌سازی در \eng{TTS}}
متن ورودی باید به شکلی تبدیل شود که سامانه بتواند آن را \textbf{درست و طبیعی بخواند}: ساختن «شکل گفتاری» (\eng{spoken form}) از «شکل نوشتاری» (\eng{written form}). انتخاب شکل گفتاری (مثلاً «شانزده و سی دقیقه» یا «چهار و نیم» و نیز «تومان» یا «تومن») به کاربرد و سبک گفتار بستگی دارد.
\end{definitionbox}

\subsection{نرمال‌سازی در \eng{TTS} با \eng{NLP} فرق دارد}
استاد تأکید کرد که نرمال‌سازی وابسته به هدف است: «متن را نسبت به چه هدفی نرمال می‌کنیم؟» نرمال‌ساز معمول \eng{NLP} (مثلاً \eng{Hazm}) \term{کاهنده}{Reductive} است: تنوع نوشتار را کم می‌کند (یکسان‌سازی «ي/ی» و «ك/ک»، نیم‌فاصله) تا برای جست‌وجو، دسته‌بندی، \eng{NER} و \eng{POS} مفید باشد. اما در \eng{TTS} متن باید \textbf{خواندنی} شود: اطلاعات نباید دور ریخته شود (مثلاً حذف \term{واژه‌های ایست}{Stopword} نامناسب است) و حتی باید اطلاعات تازه اضافه شود. روی تخته «کاهنده» در برابر «گسترش‌دهنده و آواشناسی‌محور» (\eng{Expansive \& Phonetic}) مقابل هم قرار گرفت.

\begin{warningbox}{دو نرمال‌ساز با یک نام}
نرمال‌ساز \eng{Hazm} \textbf{کاهنده} است (تنوع نوشتار را برای فهم معنا کم می‌کند)؛ نرمال‌ساز \eng{TTS} \textbf{گسترش‌دهنده} است (عدد را به کلمه تبدیل و اطلاعات تلفظی را روشن می‌کند). این دو را اشتباه نگیرید. (البته \eng{TTS} هم به یکدست‌سازی املایی برای جست‌وجو در فرهنگ تلفظ نیاز دارد، و نباید اطلاعات مؤثر بر خواندن مثل ویرگول و علامت سؤال را حذف کند.)
\end{warningbox}



\subsection{کلمات غیراستاندارد (\eng{NSW}) و روش کلاسیک}
هستهٔ نرمال‌سازی \eng{TTS} برخورد با \term{کلمات غیراستاندارد}{Non-Standard Words, NSW} است: توکن‌هایی که تلفظشان از نوشتارشان با قواعد معمول به دست نمی‌آید. دسته‌های اصلی: اعداد، تاریخ، ساعت، واحد پول، درصد، اختصارات، شماره‌ها و شناسه‌ها، و شکل‌های غیررسمی.

یک رشتهٔ ارقام را نمی‌توان با قاعدهٔ ثابتی به کلمه تبدیل کرد؛ ابتدا باید \textbf{نوع و نقش} آن مشخص شود. مثلاً در انگلیسی ۱۹۷۶ اگر سال باشد «نوزده، هفتاد و شش» خوانده می‌شود ولی در یک شناسه ممکن است رقم‌به‌رقم. در فارسی هم «/» گاهی جداکنندهٔ تاریخ ($1402/10/15$) و گاهی ممیز اعشاری ($2/5$) است و تنها نشانهٔ تفکیک بافت است (واژهٔ «تن» پس از $2/5$).

\paragraph{روش کلاسیک دو بخشی.}
\begin{enumerate}
  \item \textbf{طبقه‌بندی توکن} (\eng{Token Classification}): هر توکن به یک دسته (عدد، تاریخ، ساعت، پول، ...) نسبت داده می‌شود؛ با عبارت باقاعده (\eng{Regex})، قواعد دستی یا مدل یادگیری ماشین.
  \item \textbf{مبدل} (\eng{Converter}): مبدل مخصوص آن دسته توکن را به شکل گفتاری درمی‌آورد.
\end{enumerate}
\[
\text{\lr{Input Token}} \to \text{\lr{Token Classification}} \to \text{\lr{Category-Specific Converter}} \to \text{\lr{Spoken Form}}
\]
در کلاس کتابخانهٔ \eng{num2fawords} (تبدیل عدد به حروف فارسی؛ مثلاً words(170) برابر «یکصد و هفتاد») معرفی شد. استاد تأکید کرد که فقط برای دستهٔ \textbf{اعداد} است و تاریخ، ساعت و اختصار مبدل جدا می‌خواهند. اگر برای دسته‌ای ابزار آماده نباشد، دادهٔ آموزشی تهیه و طبقه‌بند ساخته می‌شود.

\subsection{اهمیت دامنهٔ کاربرد}
فراوانی \eng{NSW}ها به نوع متن بستگی دارد. در کلاس سه نوع داده مقایسه شد: متن خبری (\eng{New York Times})، \eng{Twitter} و \eng{SMS}.

\boardfigure{0.62}{ch07-oov-corpora.png}{احتمال دیدن پیامی با نسبت معین از توکن‌های \eng{OOV} در سه پیکره. در متن خبری حدود ۴۸ درصد پیام‌ها هیچ \eng{OOV} ندارند؛ در توییتر حدود ۳۱ درصد و در \eng{SMS} حدود ۱۵ درصد. (مقدارهای تقریبی)}

هرچه به ارتباطات غیررسمی (توییتر و \eng{SMS}) نزدیک‌تر شویم، \eng{NSW}ها (اختصار، املای غیررسمی، عامیانه) شایع‌تر می‌شوند. اگر شکل خواندن عبارتی تشخیص داده نشود، گفتار غیرطبیعی یا غلط می‌شود. این مزیت خط‌لوله است: \textbf{اهمیت و پیچیدگی هر ماژول به دامنه وابسته است}؛ برای پیام‌های کوتاه و شبکه‌های اجتماعی باید روی نرمال‌سازی بیشتر سرمایه‌گذاری کرد، برای متن رسمی می‌توان آن را ساده‌تر گرفت.

\begin{correctionbox}{\eng{OOV} و \eng{NSW} یکی نیستند}
نمودار \eng{OOV} (خارج از واژگان) را می‌شمارد؛ \eng{NSW} توکنی است که تلفظش مستقیماً از نوشتارش نمی‌آید. هم‌پوشانی زیادی دارند ولی مثلاً اشتباه تایپی \eng{OOV} است نه \eng{NSW}. در کلاس \eng{OOV} شاخصی برای فراوانی \eng{NSW} بود.
\end{correctionbox}

\section{مرحلهٔ دوم: رفع ابهام هم‌نویسه‌ها}

\begin{definitionbox}{هم‌نویسه (\eng{Homograph})}
دو یا چند واژه با نویسهٔ (املای) یکسان ولی تلفظ متفاوت. هدف \term{رفع ابهام هم‌نویسه}{Homograph Disambiguation} انتخاب تلفظ درست با توجه به بافت جمله است.
\end{definitionbox}

ریشهٔ بسیاری از هم‌نویسه‌های فارسی این است که مصوت‌های کوتاه (\eng{a}، \eng{e}، \eng{o}) در نوشتار معمولی نمایش داده نمی‌شوند:
\begin{center}
\begin{tabular}{lll}
\toprule
نوشتار & تلفظ‌های ممکن (\eng{IPA}) & معنی \\
\midrule
مرد & \eng{[mæɾd]} / \eng{[moɾd]} & انسان مذکر (اسم) / فوت کرد (فعل) \\
کرم & \eng{[keɾm]} / \eng{[keɾem]} & کرم خاکی / کرم مرطوب‌کننده \\
سرم & \eng{[seɾom]} / \eng{[sæɾæm]} & محلول پزشکی / «سرِ من» \\
\bottomrule
\end{tabular}
\end{center}

\paragraph{دو راه برای رفع ابهام.}
\begin{enumerate}
  \item \textbf{اطلاعات نحوی (\eng{Syntactic Information}):} نقش دستوری تلفظ را مشخص می‌کند. ابزار آن \term{برچسب‌گذاری اجزای کلام}{Part-of-Speech Tagging} (\eng{POS Tagging}) است که برای هر کلمه اسم، فعل، صفت، قید، ضمیر و ... را تعیین می‌کند. مثلاً «مرد» در «آن مَرد رفت» اسم و در «او مُرد» فعل است.
  \item \textbf{بافت و معنا (\eng{Context}):} اگر نقش دستوری یکسان باشد (کرم خاکی و کرم مرطوب‌کننده هر دو اسم‌اند) باید معنای جمله را به کار برد؛ یک راه یادگیری ماشینی \term{نمایش بافتی}{Contextual Embedding} است که بردار هر کلمه را با توجه به جملهٔ اطراف می‌سازد (جزئیات در جلسات بعد).
\end{enumerate}

\subsection{مدلی برای رفع ابهام هم‌نویسه}
روی تخته مدلی نمونه رسم شد: جمله از یک مدل \eng{Embedding} می‌گذرد؛ نمایش کلمهٔ هدف («گرم») به یک طبقه‌بند داده می‌شود که فقط میان تلفظ‌های \textbf{مرتبط با همان کلمه} (\eng{garm} و \eng{geram}) بیشینه می‌گیرد.

\boardfigure{0.5}{ch07-homograph-model.png}{مدل نمونهٔ رفع ابهام هم‌نویسه: \eng{Embedding} بردار بافتی را می‌سازد و طبقه‌بند تنها میان تلفظ‌های نامزد کلمهٔ هدف بیشینه می‌گیرد.}

چرا یک طبقه‌بند \textbf{سراسری} (\eng{Global Classifier}) برای همهٔ هم‌نویسه‌ها، نه یکی برای هر کلمه؟ هم‌نویسه‌ها در قواعد نحوی و معنایی مشترک‌اند؛ طبقه‌بند جدا برای کلمات کم‌تکرار بیش‌برازش می‌کند یا چیزی یاد نمی‌گیرد؛ با وزن‌های مشترک (\eng{Shared Weights}) دانشی که از کلمات پرتکرار یاد گرفته شده روی کلمات کم‌تکرار هم اعمال می‌شود (\term{یادگیری انتقالی میان‌کلمه‌ای}{Cross-word Transfer Learning}).

\begin{exambox}{پرسش روی تخته: «اینجا \eng{Argmax} یا تصادفی؟»}
در رفع ابهام هم‌نویسه تلفظ \textbf{درست و قطعی} می‌خواهیم، پس معمولاً \eng{argmax} روی نامزدهای مجاز انتخاب می‌شود؛ نمونه‌برداری تصادفی بیشتر برای تنوع خروجی‌های تولیدی (مثلاً پروزودی) مطرح است.
\end{exambox}

\begin{correctionbox}{هر ابهام معنایی ابهام تلفظی نیست}
رفع ابهام هم‌نویسه به \term{رفع ابهام معنایی واژه}{Word Sense Disambiguation, WSD} نزدیک است، اما هدفش تعیین \textbf{تلفظ} است. هر سه معنای «شیر» (جانور، نوشیدنی، شیر آب) در فارسی \eng{[ʃiɾ]} تلفظ می‌شود؛ پس ابهام معنایی (\eng{WSD}) است ولی برای \eng{TTS} هم‌نویسه نیست. اطلاعات معنایی فقط ابزاری برای رسیدن به تلفظ صحیح است.
\end{correctionbox}

\section{تبدیل نویسه به واج (\eng{G2P})}

تا اینجا ورودی توالی نویسه‌ها (\eng{Grapheme}) بود؛ برای تولید گفتار به \term{واج}{Phoneme} نیاز داریم. این تبدیل را \term{نویسه‌به‌واج}{Grapheme-to-Phoneme, G2P} می‌نامند.

\begin{definitionbox}{نویسه و واج}
\textbf{نویسه (\eng{Grapheme})}: واحد تمایزدهندهٔ نظام نوشتاری (مثل «س»، \eng{a}). \\
\textbf{واج (\eng{Phoneme})}: کوچک‌ترین واحد صوتی تمایزدهنده در نظام آوایی زبان؛ با الفبای آوانگاری مانند \eng{IPA} نوشته می‌شود.
\end{definitionbox}

\begin{correctionbox}{نویسه «معنادار» نیست}
روی تخته نویسه «کوچک‌ترین واحد معنادار خط» نوشته شده بود؛ دقیق‌تر: کوچک‌ترین واحد \textbf{متمایزکنندهٔ} خط. «معنادار بودن» ویژگی \term{تکواژ}{Morpheme} است.
\end{correctionbox}

نگاشت نویسه به واج یک‌به‌یک نیست (انگلیسی هم: \eng{tion} $\leftarrow$ \eng{ʃən}). در فارسی:
\begin{itemize}
  \item مصوت‌های کوتاه نوشته نمی‌شوند ولی در نمایش واجی باید بیایند: استاد «کتاب» را \eng{ketab} و «شب» را \eng{ʃab} نوشت؛
  \item تشدید، سکون و استثناهای املایی باید در نظر گرفته شوند؛
  \item بعضی نویسه‌ها چند تلفظ دارند («و»: /v/، /u/، /o/، /ow/) و بعضی واج‌ها چند نویسه (/s/: «س»، «ص»، «ث»).
\end{itemize}

\paragraph{سه پیش‌نیاز طراحی \eng{G2P}.}
\begin{enumerate}
  \item \textbf{الفبای آوانگاری} (\term{الفبای آوایی بین‌المللی}{International Phonetic Alphabet, IPA}): اصل آن «یک صدا، یک نماد» است و جدولش بر پایهٔ همان ویژگی‌های تولیدی (محل و شیوهٔ تولید برای صامت‌ها؛ ارتفاع و پیشین/پسین بودن زبان برای مصوت‌ها) سازمان یافته است.
  \item \textbf{فرهنگ تلفظ} (\eng{Pronunciation Dictionary}): نگاشت نوشتار به تلفظ؛ هم برای یادگیری و هم برای سنجش مدل. نمونه‌ها: \eng{KaamelDict} (فرهنگ تلفظ فارسی روی \eng{Hugging Face}) و \eng{Persian Phonetic Corpus}؛ ابزار \eng{phonemizer} متن را به \eng{IPA} تبدیل می‌کند.
  \item \textbf{مدل \eng{G2P}}: مثلاً مدل دنباله‌به‌دنبالهٔ رمزگذار–رمزگشا (\eng{Encoder--Decoder}) که توالی نویسه‌ها را به توالی واج‌ها می‌برد (معماری در جلسات بعد).
\end{enumerate}

\paragraph{طرح عملی پیشنهادی.} روی تخته نمودار تصمیمی رسم شد: اگر کلمه در فرهنگ تلفظ نباشد (\eng{OOV})، یک مدل رمزگذار–رمزگشا در سطح بایت (\eng{ByT5}) که روی تلفظ کلمات فرهنگ آموزش دیده و نوشتار را بایت‌به‌بایت (\eng{UTF-8}) می‌خواند، قواعد پنهان آوایی را تعمیم می‌دهد؛ اگر در فرهنگ باشد و هم‌نویسه باشد، طبقه‌بند سراسری تلفظ را برمی‌گزیند؛ وگرنه تلفظ ثبت‌شده یا کتابخانهٔ دست‌ساز به کار می‌رود.

\begin{figure}[!htbp]
\centering
\resizebox{\textwidth}{!}{%
\begin{tikzpicture}[rtl,
  dec/.style={draw=mainblue, thick, diamond, aspect=2.4, fill=softyellow, align=center, inner sep=2pt, font=\normalsize},
  box/.style={draw=mainblue, thick, rounded corners=3pt, fill=softblue, align=center, minimum height=1.1cm, minimum width=3.2cm, font=\normalsize},
  arr/.style={-{Stealth[length=2.8mm]}, thick, mainblue},
  lab/.style={font=\small, fill=white, inner sep=1.5pt},
  note/.style={draw=black!45, dashed, rounded corners=2pt, align=center, font=\footnotesize, text=black!75, inner sep=3pt}
]
\node[dec] (oov) at (4.6,0) {\lr{OOV?}};
\node[dec] (hom) at (9.4,-2.6) {هم‌نویسه؟};
\node[box, fill=softgreen, minimum width=4.2cm] (byt5) at (0.4,-2.6) {\lr{ByT5}\\\lr{Encoder--Decoder}};
\node[box] (glob) at (7.0,-5.2) {\lr{Global Classifier}};
\node[box] (lib) at (11.6,-5.2) {کتابخانه‌ها /\\دست‌ساز};
\draw[arr] (oov.west) -- node[lab, above left] {بله} (byt5.north);
\draw[arr] (oov.east) -- node[lab, above right] {خیر} (hom.north);
\draw[arr] (hom.south west) -- node[lab, above left] {بله} (glob.north);
\draw[arr] (hom.south east) -- node[lab, above right] {خیر} (lib.north);
\node[note] (n1) at (0.4,-4.7) {\rl{آموزش روی تلفظ واژه‌های فرهنگ}\\\rl{خواندن در سطح بایت‌های \lr{UTF-8}}\\\rl{حالت نهایی رمزگذار (\lr{Context Vector}) به رمزگشا می‌رود}};
\draw[dashed, black!45] (byt5.south) -- (n1.north);
\node[note] (q1) at (-0.6,-0.1) {\rl{پرسش روی تخته:}\\\rl{«تسک چیست؟»}};
\draw[dashed, black!45] (q1.south) -- (byt5.north west);
\node[note] (q2) at (7.0,-6.9) {\rl{پرسش روی تخته:}\\\rl{«اینجا \lr{Argmax} یا نمونه‌گیری تصادفی؟»}};
\draw[dashed, black!45] (glob.south) -- (q2.north);
\end{tikzpicture}}
\caption{نمودار تصمیم برای بخش تلفظ: کلمهٔ \eng{OOV} به \eng{ByT5}، کلمهٔ هم‌نویسه به طبقه‌بند سراسری و در غیر این صورت به فرهنگ/کتابخانهٔ دست‌ساز می‌رود. (بازسازی‌شده از تختهٔ کلاس)}
\end{figure}

\begin{exambox}{پایان بخش تحلیل متن}
در پایان \eng{G2P} خروجی یک \textbf{توالی واجی} است؛ در این نقطه \eng{Text Analysis} (\eng{Front-End}) پایان می‌یابد و تولید صوت (\eng{Back-End}) آغاز می‌شود.
\end{exambox}

\section{مرحلهٔ دوم سامانه: تولید گفتار از روی واج‌ها}

ورودی مجموعه‌ای از واج‌هاست و خروجی موج گفتار. سه چالش اصلی داریم:
\begin{enumerate}
  \item \textbf{اطلاعات پروزودی:} متن لحن و آهنگ را مشخص نمی‌کند؛ «سلام» می‌تواند عادی، پرسشی، متعجبانه یا با تأکید گفته شود. بخشی از اطلاعات از نشانه‌های نگارشی و بافت به دست می‌آید، ولی کافی نیست؛ پس مدل باید پروزودی را استنتاج کند یا کنترل‌های جدا (نوع صدا، سبک، حالت) بگیرد تا گفتار رباتی نباشد.
  \item \textbf{اختلاف نرخ و ابعاد:} ۳۰ واج ورودی در برابر موج ۳ ثانیه‌ای در ۲۴ کیلوهرتز:
  \[
  N = f_s \times T = 24000 \times 3 = 72\,000 .
  \]
  اگر طیف‌نگار مل ۸۰ بعدی با گام ۳۰۰ نمونه داشته باشیم، میانگین هر واج ($75$ میلی‌ثانیه) حدود ۶ قاب و ۱۸۰۰ نمونه است. در کلاس به این افزایش وضوح زمانی \eng{Upsampling} گفته شد؛ اما مسئله فقط زیادکردن تعداد نمونه نیست و مدل باید مقدار هر نمونه را هم تعیین کند.
  \item \textbf{تطابق زمانی (\eng{Temporal Alignment}):} هر واج کی شروع شود، چقدر طول بکشد، مرز دو واج کجاست و کدام واج‌ها (به‌دلیل هم‌تولیدی) بر هم اثر می‌گذارند. این به \term{مدل‌سازی مدت}{Duration Modeling} مربوط است؛ واج‌ها مدت یکسان ندارند.
\end{enumerate}

\begin{correctionbox}{نرخ ۲۴ کیلوهرتز و واژهٔ «\eng{Upsampling}»}
نرخ نمونه‌برداری «استاندارد» واحدی وجود ندارد (تلفن ۸، تشخیص گفتار معمولاً ۱۶، پژوهش‌های \eng{TTS} عمدتاً $22.05$ یا ۲۴ کیلوهرتز)؛ ۲۴ کیلوهرتز مثال معقولی است. همچنین «\eng{Upsampling}» اصطلاحی تقریبی است: افزایش طول در دو جا رخ می‌دهد، واج به قاب (با مدل‌سازی مدت) و قاب به نمونه (چند صد نمونه به‌ازای هر قاب).
\end{correctionbox}

\section{مدل آکوستیکی}

برای ساده‌کردن مسئله، تولید گفتار به دو مرحله تقسیم می‌شود:
\[
\text{\lr{Phonemes}} \xrightarrow{\;\text{\lr{Acoustic Model}}\;} \text{\lr{Mel-Spectrogram}} \xrightarrow{\;\text{\lr{Vocoder}}\;} \text{\lr{Waveform}}
\]
\term{مدل آکوستیکی}{Acoustic Model} از روی واج‌ها نمایش آکوستیکی (مثل طیف‌نگار مل) می‌سازد و باید زمان‌بندی، ویژگی‌های آکوستیکی و در صورت لزوم پروزودی را لحاظ کند.

\paragraph{چرا طیف‌نگار به‌عنوان نمایش میانی؟} پاسخ کلاس: نمایش میانی \textbf{پیچیدگی مسئلهٔ یادگیری را کاهش می‌دهد}؛ به‌جای ساختن همهٔ جزئیات موج در یک گام، ابتدا نمایشی ساختاریافته‌تر ساخته می‌شود و ماژول دیگری آن را به موج برمی‌گرداند.

\begin{exambox}{مثال استاد: مسیر تهران تا مشهد}
فرض کنید فقط می‌دانیم کسی از تهران به مشهد رفته و مسافت مشخصی را طی کرده است؛ بازسازی مسیر دقیقش دشوار است (از کدام شهرها و جاده‌ها گذشته؟). اگر «نقاط عطف» میان راه را بدانیم، بازسازی ساده‌تر می‌شود. در تولید گفتار هم دانستن واج‌ها برای تعیین همهٔ جزئیات موج کافی نیست و طیف‌نگار نقش همان نقاط عطف را دارد.
\end{exambox}

\paragraph{نقش \eng{Embedding}.} برای هر واج یک بردار \term{تعبیه}{Embedding} (مثلاً ۵۱۲ بعدی؛ عدد فقط مثال است) یاد گرفته می‌شود که اطلاعات واج و ارتباطش با واج‌های قبل و بعد را حمل می‌کند؛ پس مدل فقط به هویت یک واج مستقل محدود نیست و می‌تواند هم‌تولیدی را بیاموزد.

\paragraph{مدل آکوستیکی و پروزودی.} استاد تأکید کرد که «تولید طیف‌نگار» و «پیش‌بینی پروزودی» دقیقاً یک مسئله نیستند (برای یک محتوای ثابت حالت‌های پروزودی گوناگون ممکن است)؛ بنابراین در بعضی معماری‌ها برای پروزودی و خروجی آکوستیکی مؤلفه‌ها و تابع هزینه‌های جدا تعریف می‌شود (مثلاً \eng{FastSpeech 2}: پیش‌بین‌های جدا برای مدت، گام و انرژی).

\section{وکودر: از طیف‌نگار به موج}

ماژولی که نمایش آکوستیکی را به \eng{Waveform} تبدیل می‌کند \term{وکودر}{Vocoder} نام دارد (از «\eng{Voice + coder}»).

\subsection{چرا ساختن موج از طیف‌نگار ساده نیست؟}
ضرایب تبدیل فوریه هم \textbf{دامنه} دارند و هم \textbf{فاز}. برای ضریب $X = a + jb$:
\[
|X| = \sqrt{a^{2}+b^{2}}, \qquad \phi = \text{\lr{atan2}}(b,a).
\]
طیف‌نگار فقط اندازه (\eng{Magnitude}) را نگه می‌دارد و فاز را (که مشخص می‌کند هر مؤلفهٔ فرکانسی کجای زمان است و چگونه با بقیه ترکیب می‌شود) کنار می‌گذارد. پس بدون اطلاعات یا تخمین اضافی نمی‌توان همهٔ جزئیات موج اولیه را بازیافت.

\begin{addedbox}{مثال عددی: اندازه یکسان، فاز متفاوت}
$X_1 = 3+4j$ و $X_2 = -3+4j$ هر دو اندازهٔ $\sqrt{9+16}=5$ دارند، ولی $\phi_1 \approx 53.1^\circ$ و $\phi_2 \approx 126.9^\circ$. از روی اندازهٔ ۵ نمی‌توان فهمید ضریب اصلی کدام بوده است. شکل زیر همین را روی سیگنال نشان می‌دهد: دو سیگنال با طیف اندازهٔ یکسان و شکل موج کاملاً متفاوت.
\end{addedbox}

\coursefigure{0.6}{ch07-phase-demo.png}{دو سیگنال با طیف اندازهٔ یکسان و فازهای متفاوت: ردیف بالا شکل موج و ردیف پایین طیف اندازه. (شکل افزوده)}

\subsection{بازسازی فاز و مدل‌های مولد}
برای ساختن موج باید فاز را بازسازی یا سیگنالی سازگار تولید کرد:
\begin{itemize}
  \item \textbf{کلاسیک (\eng{Griffin--Lim}):} فاز را تکراری حدس می‌زند تا طیف‌نگار سیگنال با طیف‌نگار مفروض سازگار شود؛ ساده و بدون آموزش، اما صدایی مصنوعی و «وزوزی» دارد.
  \item \textbf{وکودرهای عصبی:} شبکه مستقیماً از طیف‌نگار به موج می‌رسد؛ مثلاً \eng{WaveNet} (خودهمبسته، نمونه‌به‌نمونه و کند) و \eng{HiFi-GAN} (سریع، مبتنی بر \eng{GAN}).
\end{itemize}

\begin{correctionbox}{فاز «کاملاً» از بین نمی‌رود}
چون پنجره‌های طیف‌نگار هم‌پوشان‌اند، افزونگی وجود دارد و فاز را می‌توان از سازگاری میان قاب‌ها تا حد خوبی تخمین زد (کاری که \eng{Griffin--Lim} می‌کند). طیف‌نگار مل علاوه بر فاز تفکیک فرکانسی را هم کم می‌کند؛ پس ساخت موج از مل مسئله‌ای تقریباً بدوضع‌طرح‌شده است که با مدل مولد خوب حل می‌شود.
\end{correctionbox}

\paragraph{\eng{GAN}.} \term{شبکهٔ مولد رقابتی}{Generative Adversarial Network, GAN} دو بخش دارد: \term{مولد}{Generator} که خروجی مصنوعی (اینجا موج) می‌سازد و \term{تمایزدهنده}{Discriminator} که می‌کوشد آن را از نمونهٔ واقعی تشخیص دهد؛ بازخورد تمایزدهنده مولد را بهتر می‌کند. در پاسخ به پرسشی در کلاس، استاد روشن کرد که مدل لازم نیست «فاز» مستقل تولید کند؛ می‌تواند مستقیماً موج بسازد و با تابع‌های هزینه کیفیتش را بهتر کند. جزئیات به جلسات بعد موکول شد.

\section{جمع‌بندی خط‌لولهٔ \eng{TTS}}

\begin{center}
\small
\begin{tabular}{llll}
\toprule
مرحله & ورودی & خروجی & وظیفه \\
\midrule
\eng{Text Normalization} & متن خام & متن قابل خواندن & گسترش \eng{NSW} \\
\eng{Homograph Disamb.} & متن دارای هم‌نویسه & تلفظ مشخص‌شده & رفع ابهام بافتی و نحوی \\
\eng{G2P} & نویسه‌ها & واج‌ها & تبدیل نوشتار به نمایش واجی \\
\eng{Acoustic Model} & توالی واج‌ها & طیف‌نگار (مل) & تولید نمایش آکوستیکی \\
\eng{Vocoder} & طیف‌نگار & موج & تولید موج گفتار \\
\bottomrule
\end{tabular}
\end{center}

\paragraph{چرا خط‌لوله هنوز مهم است؟} کنترل مستقل هر بخش (تنظیم نرمال‌ساز برای یک دامنه، کاهش تأخیر، مدل کلاسیک برای بخشی از سامانه، تمرکز روی اعداد یا هم‌نویسه‌ها) مهم است؛ انتخاب معماری باید بر پایهٔ مسئله و محدودیت‌های کاربرد باشد.

\begin{exambox}{ارزیابی و جلسات آینده}
در پایان کلاس قرار شد مفاهیمی مانند \eng{Transformer}، \eng{Self-Attention}، \eng{Cross-Attention}، \eng{Encoder--Decoder} و \eng{GAN} به‌اندازهٔ لازم توضیح داده شوند تا بتوان به آن‌ها ارجاع داد؛ یک جلسه هم به \eng{End-to-End} اختصاص می‌یابد. استاد اشاره کرد که مطالب تدریس‌شده می‌توانند در امتحان بیایند. تکلیف مشخصی اعلام نشد.
\end{exambox}

\section{جمع‌بندی}
در این جلسه از تحلیل سیگنال به \textbf{تولید} سیگنال رسیدیم. \eng{TTS} دشوار است چون متن بسیار کم‌حجم‌تر از موج است و اطلاعات پروزودی و گوینده را ندارد. سامانه به \eng{Front-End} (نرمال‌سازی، رفع ابهام هم‌نویسه، \eng{G2P}) و \eng{Back-End} (مدل آکوستیکی و وکودر) تقسیم می‌شود. در بخش اول «خواندنی کردن» متن مهم است و با نرمال‌سازی \eng{NLP} فرق دارد؛ در بخش دوم طیف‌نگار مل پیچیدگی را کم می‌کند و وکودر باید فاز ازدست‌رفته را بازسازی کند.

\begin{exambox}{نکات مهم برای مرور و امتحان}
\begin{enumerate}
  \item \eng{TTS}: متن $\leftarrow$ موج. چالش‌ها: شکاف حجم اطلاعات، کم‌گویی متن، نگاشت یک‌به‌چند.
  \item دو رویکرد \eng{Pipeline} و \eng{End-to-End}؛ سرتاسری لزوماً بهتر نیست (تأخیر، دامنه، کنترل).
  \item \eng{Front-End} (متن $\leftarrow$ واج) و \eng{Back-End} (واج $\leftarrow$ طیف‌نگار $\leftarrow$ موج).
  \item نرمال‌سازی \eng{TTS} گسترش‌دهنده و آواشناسی‌محور است، نه کاهنده مانند \eng{Hazm}.
  \item \eng{NSW}: ابتدا تشخیص دستهٔ توکن، سپس مبدل همان دسته؛ فراوانی \eng{NSW} به دامنه وابسته است (توییتر و \eng{SMS} بیشتر از خبر).
  \item هم‌نویسه‌ها با \eng{POS} و بافت (\eng{Contextual Embedding}) حل می‌شوند؛ هر ابهام معنایی ابهام تلفظی نیست.
  \item \eng{G2P}: سه پیش‌نیاز: الفبای آوایی (\eng{IPA})، فرهنگ تلفظ، مدل؛ نگاشت یک‌به‌یک نیست.
  \item چالش‌های \eng{Back-End}: پروزودی، اختلاف نرخ (۳۰ واج در برابر $72\,000$ نمونه)، هم‌ترازی زمانی و هم‌تولیدی.
  \item مدل آکوستیکی واج را به طیف‌نگار مل می‌برد؛ نمایش میانی پیچیدگی را کم می‌کند (مثال تهران–مشهد).
  \item وکودر طیف‌نگار را به موج می‌برد؛ مشکل اصلی: از دست رفتن فاز. راه‌حل‌ها: \eng{Griffin--Lim} و وکودرهای عصبی (\eng{WaveNet}، \eng{HiFi-GAN}).
\end{enumerate}
\end{exambox}
